% Generated draft: requires mathematical and source review.
\documentclass[11pt]{article}
\usepackage[margin=25mm]{geometry}
\usepackage{fontspec}
\setmainfont{Latin Modern Roman}
\usepackage{amsmath,amssymb,mathtools,unicode-math}
\setmathfont{Latin Modern Math}
\usepackage{enumitem,needspace,xurl,hyperref}
\setlength{\parindent}{0pt}
\setlength{\parskip}{5pt}
\setlength{\emergencystretch}{3em}
\newcommand{\N}{\mathbb{N}}
\newcommand{\Z}{\mathbb{Z}}
\newcommand{\Q}{\mathbb{Q}}
\newcommand{\R}{\mathbb{R}}
\newcommand{\C}{\mathbb{C}}
\newcommand{\F}{\mathbb{F}}
\newcommand{\A}{\mathbb{A}}
\newcommand{\PP}{\mathbb{P}}
\newcommand{\E}{\mathbb{E}}
\newcommand{\eps}{\varepsilon}
\newcommand{\dd}{\,\mathrm{d}}
\begin{document}
\section*{UnsolvedMath ID: 11000320}
\textbf{Source problem code:} catalogue entry
\subsection{Definitions and mathematical statement}
Problem 8.1 asks to determine explicitly the image and cokernel of the homomorphism labeled (15). Its formula, domain, codomain, grading and coefficients are absent. The surrounding text places it in the Johnson-homomorphism setting with target notation \(\mathfrak h_{g,1}\), but does not define that object here.

For a specified additive homomorphism \(\phi:A\to B\), the targets would be \(\operatorname{im}\phi\) and \(\operatorname{coker}\phi=B/\operatorname{im}\phi\). If the source is using graded modules or Lie algebras, the category must be recovered: a Lie-subalgebra image is not necessarily an ideal, so a vector-space cokernel and a Lie-algebra quotient are not automatically identical. No particular Johnson filtration or integral lattice is silently supplied.

The source reports Hain [29] proving that, after tensoring with \(\Q\), the image of (15) is precisely the Lie subalgebra generated by the source's degree-one part. It says identifying which elements of \(\mathfrak h_{g,1}^{\Q}\) lie in that subalgebra remains unclear. The rational structural description is not a fully explicit image/cokernel calculation and does not determine integral torsion.

The context recounts Oda's late-1980s prediction of arithmetic obstructions, motivated by Ihara's work mainly in parameters \(g=0,n=3\). It says the absolute Galois group \(\operatorname{Gal}(\overline\Q/\Q)\) should ``appear'' in
\[
 \mathfrak h_{g,1}\otimes\Z_\ell
\]
outside the geometric part, in an \(\mathrm{Sp}\)-invariant way, for every genus \(g\) and prime \(\ell\). The original absolute-Galois phrase determines the algebraic-closure reading despite flattened bars. The relevant quotient, graded/profinite realization, geometric part and invariance are undefined, so ``appear'' is not promoted to a literal embedding of the full Galois group into an additive Lie algebra.

The source credits Nakamura [86] in 1994, with Matsumoto [65] also cited, for results establishing the predicted phenomenon in its precise source sense. It calls this a second Johnson-surjectivity obstruction after the traces in [76]. The exact theorem is absent, and the final sentence introduces a following unprovided problem; neither is invented.
\subsection{Short English statement}
Compute the source's Johnson-type homomorphism image and cokernel explicitly. The map itself is missing. Hain's rational degree-one-generation report and the historical arithmetic-obstruction discussion are retained with their integral, categorical and Galois-realization qualifications.
\subsection{Sources}
\begin{itemize}
\item Local UnsolvedMath registry, permanent ID \(11000320\), exact source code \texttt{catalogue entry}, Problem 8.1, PDF page 373. No individual proposer or proposed year is supplied; the title is truncated.
\item Farb, edited volume \emph{Problems on Mapping Class Groups and Related Topics}: \url{https://www.math.uchicago.edu/~farb/papers/mcgbook.pdf}. The volume has 329 explicitly typeset numbered items across chapters, repeated local numbering distinguished by PDF page. The extract contains chapter/page furniture on cohomological structure and page 367.
\item Internal references: Hain [29]; Nakamura [86] (1994); Matsumoto [65]; traces [76]. Oda's late-1980s prediction and Ihara's number-theoretic work are named without separate complete citations. Homomorphism (15) and the next problem are not supplied.
\item Extraction pdf-font-structure, accessed 2026-07-29. Source evidence NEEDS\_REVIEW presents an open problem and requires a current primary-literature audit. Registry status open and classification OPEN-TRIAGE; research summary explicitly cannot verify status.
\item Dataset attribution: Ulam AI / UnsolvedMath Contributors, CC BY 4.0. Source rights NEEDS\_REVIEW: publicly accessible author/publisher copy, redistribution terms requiring review.
\end{itemize}
\end{document}
